\documentclass{article}
\usepackage[a4paper, total={6.2in, 10in}]{geometry}
\usepackage{enumitem}
\usepackage{amssymb}
\usepackage{bussproofs}

\title{Práctica 3 - Demostración en Lógica Proposicional}
\author{Paradigmas de Programación}
\date{1C 2025}

\begin{document}
\renewcommand{\baselinestretch}{1.5} 
\maketitle

\section*{Ejercicio 1}{Determinar el valor de verdad de las siguientes proposiciones (fórmulas) cuando el valor de verdad de P y Q es \textbf{V}, mientras que el de S y T es \textbf{F}.}
{\begin{enumerate}[label=\roman*.,font=\itshape]
\item {$(\lnot P \lor Q) \rightarrow \textbf{V}$}
\item {$(P \lor (S \land T) \lor Q) \rightarrow \textbf{V}$}
\item {$\lnot (Q \lor S) \rightarrow \textbf{F}$}
\item {$(\lnot Q \lor S) \iff (\lnot P \land \lnot S) \rightarrow \textbf{V}$}
\item {$((P \lor S) \land (T \lor Q)) \rightarrow \textbf{V}$}
\item {$(((P \lor S) \land (T \lor Q)) \iff (P \lor (S \land T ) \lor Q)) \rightarrow \textbf{V}$}
\item {$(\lnot Q \land \lnot S) \rightarrow \textbf{F}$}
\end{enumerate}}
\section*{Ejercicio 2}{Mostrar que cualquier fórmula de la lógica proposicional que utilice los conectivos {$\lnot$} (negación), {$\land$} (conjunción),
{$\lor$} (disyunción), {$\Rightarrow$} (implicación) puede reescribirse a otra fórmula equivalente que usa sólo los conectivos {$\lnot$} y {$\lor$}. Sugerencia: hacer inducción en la estructura de la fórmula.}
\\
\\
{Idea general: hacer inducción en la estructura de los conectivos lógicos. Tengo una fórmula $\tau$ que puede usar el conjunto de conectivos lógicos \{ $\lnot, \ \land, \ \lor, \ \Rightarrow$\} y la quiero reescribir como una fórmula $\sigma$ que sólo utilice el conjunto de conectivos lógicos \{$\lnot, \ \lor$\} }.
\\
\\
{\textbf{Caso base}: fórmula atómica. \\ $\tau_{1}$ no usa ningún conectivo lógico. $\tau_{1}$ equivale a $\sigma_{1}$ que tampoco usa ningún conectivo lógico.}
\\
\\
{\textbf{Caso $\lnot$}: \\ $\tau_{1}$ sólo utiliza el conectivo $\lnot$ y es equivalente a $\sigma_{1}$ que también cumple.}
\\
\\
{\textbf{Caso $\lor$}: \\ $\tau_{1} = \sigma_{1} \lor \sigma_{2}$ es equivalente a $\tau_{2} = \sigma_{1} \lor \sigma_{2}$ es equivalente a $\tau_2$ que sólo utiliza el conectivo $\lor$.} 
\\
\\
{\textbf{Caso $\Rightarrow$}: \\ $\tau_{1} = \sigma_{1}\Rightarrow\sigma_{2}$ es equivalente a $\tau_{2} = \lnot\sigma_{1} \lor \sigma_{2}$ que sólo utiliza el conectivo $\lor$.}
\\
\\
{\textbf{Caso $\land$}: \\ $\tau_{1} = \sigma_{1} \land \sigma_{2}$ es equivalente a $\tau_{2}=\lnot(\lnot\sigma_{1} \lor \lnot\sigma_{2}$ que sólo utiliza los conectivos $\lnot$ y $\lor$.}
\section*{Ejercicio 3}{Sean $\tau, \ \sigma, \ \rho, \ $ y $\zeta$ proposiciones tales que $\tau \Rightarrow \sigma$ es tautología y $\rho \Rightarrow \zeta$ es contradicción. Determinar si las siguientes proposiciones son tautologías, contradicciones o contingencias y demostrarlo:}
\begin{enumerate}[label=\roman*.,font=\itshape]
\item{$(\tau \Rightarrow \sigma) \lor (\rho \Rightarrow \zeta)$}
\\
\\
{V $\vDash (\tau \Rightarrow \sigma) \lor (\rho \Rightarrow \zeta) \iff $ V $\vDash (\tau \Rightarrow \sigma)$ ó V $\vDash (\rho \Rightarrow \zeta)$ . Importante: la V es de Valuación!!}
\\
\\
{Como sé que $(\tau \Rightarrow \sigma)$ es una tautología, entonces V $\vDash (\tau \Rightarrow \sigma)$ para cualquier V. Luego,\\ V $\vDash (\tau \Rightarrow \sigma) \lor (\rho \Rightarrow \zeta) $ para cualquier V.}
\end{enumerate}

\section*{Ejercicio 5}{Demostrar en deducción natural que las siguientes fórmulas son teoremas \textbf{sin usar principios de razonamiento clásicos} salvo que se indique lo contrario. Recordemos que una fórmula $\tau$ es un teorema si y sólo si vale $\vdash \sigma$:}
\begin{enumerate}[label=\roman*.,font=\itshape]
\item{$Modus \ ponens$ relativizado:
$(\rho \Rightarrow \sigma \Rightarrow \tau ) \Rightarrow (\rho \Rightarrow \sigma) \Rightarrow \rho \Rightarrow \tau$}
\begin{prooftree}
            \AxiomC{}
            \RightLabel{$_{Ax.}$}
            \UnaryInfC{$\Gamma \vdash \rho \Rightarrow \sigma \Rightarrow \tau$}
            
            \AxiomC{}
            \RightLabel{$_{Ax.}$}
            \UnaryInfC{$\Gamma \vdash \rho$}
        \RightLabel{$\Rightarrow_e$}
        \BinaryInfC{$\Gamma \vdash \sigma \Rightarrow \tau$}

            \AxiomC{}
            \RightLabel{$_{Ax.}$}
            \UnaryInfC{$\Gamma \vdash \rho \Rightarrow \ \sigma$}

            \AxiomC{}
            \RightLabel{$_{Ax.}$}
            \UnaryInfC{$\Gamma \vdash \rho$}
        \RightLabel{$\Rightarrow_e$}
        \BinaryInfC{$\Gamma \vdash \sigma$}
    
    \RightLabel{$\Rightarrow_e$}
    \BinaryInfC{$\Gamma = (\rho \Rightarrow \sigma \Rightarrow \tau ), \ (\rho \Rightarrow \sigma), \ \rho \ \vdash \tau$}
    
    \RightLabel{$\Rightarrow_i$}
    \UnaryInfC{$(\rho \Rightarrow \sigma \Rightarrow \tau ), \ (\rho \Rightarrow \sigma) \vdash \rho \Rightarrow \tau$}
    
    \RightLabel{$\Rightarrow_i$}
    \UnaryInfC{$(\rho \Rightarrow \sigma \Rightarrow \tau ) \vdash (\rho \Rightarrow \sigma) \Rightarrow \rho \Rightarrow \tau$}
    
    \RightLabel{$\Rightarrow_i$}
    \UnaryInfC{$\vdash (\rho \Rightarrow \sigma \Rightarrow \tau ) \Rightarrow (\rho \Rightarrow \sigma) \Rightarrow \rho \Rightarrow \tau$}
\end{prooftree}
\item{Reducción al absurdo: $(\rho \Rightarrow \bot) \Rightarrow \lnot \rho$}
\begin{prooftree}
        \RightLabel{$_{Ax.}$}
        \AxiomC{}
        \UnaryInfC{$(\rho \Rightarrow \bot), \ \rho \ \vdash \rho \Rightarrow \bot$}
        
        \RightLabel{$_{Ax.}$}
        \AxiomC{}
        \UnaryInfC{$(\rho \Rightarrow \bot), \ \rho \ \vdash \rho$}

    \RightLabel{$\Rightarrow_e$}
    \BinaryInfC{$(\rho \Rightarrow \bot), \ \rho \ \vdash \bot$}
        
    \RightLabel{$\lnot_i$}
    \UnaryInfC{$(\rho \Rightarrow \bot) \vdash \lnot \rho$}
    
    \RightLabel{$\Rightarrow_i$}
    \UnaryInfC{$\vdash (\rho \Rightarrow \bot) \Rightarrow \lnot \rho$}
\end{prooftree}

\item {Introducción de la doble negación: $\rho \Rightarrow \lnot\lnot\rho$}
\begin{prooftree}
        \RightLabel{$_{Ax.}$}
        \AxiomC{}
        \UnaryInfC{$\rho, \ \rho \ \vdash \rho$}
        
        \RightLabel{$_{Ax.}$}
        \AxiomC{}
        \UnaryInfC{$\rho, \ \rho \ \vdash \lnot \rho$}

    \RightLabel{$\lnot_e$}
    \BinaryInfC{$\rho, \ \lnot\rho \ \vdash \bot$}

    \RightLabel{$\lnot_i$}
    \UnaryInfC{$\rho \ \vdash \lnot\lnot\rho$}
    
    \RightLabel{$\Rightarrow_i$}
    \UnaryInfC{$\vdash \rho \Rightarrow \lnot\lnot\rho$}
\end{prooftree}

\item {Eliminación de la triple negación: $\lnot\lnot\lnot\rho \Rightarrow \lnot\rho$}
\begin{prooftree}
        \RightLabel{$_{Ax.}$}
        \AxiomC{}
        \UnaryInfC{$\lnot\lnot\lnot\rho, \ \rho \ \vdash \lnot\lnot\lnot\rho$}

            \RightLabel{$_{Ax.}$}
            \AxiomC{}
            \UnaryInfC{$\lnot\lnot\lnot\rho, \ \lnot\rho, \ \rho \ \vdash \lnot\rho$}
            
            \RightLabel{$_{Ax.}$}
            \AxiomC{}
            \UnaryInfC{$\lnot\lnot\lnot\rho, \ \lnot\rho, \ \rho \ \vdash \rho$}
            
        \RightLabel{$\lnot_e$}
        \BinaryInfC{$\lnot\lnot\lnot\rho, \ \lnot\rho, \ \rho \ \vdash \bot$}

        \RightLabel{$\lnot_i$}
        \UnaryInfC{$\lnot\lnot\lnot\rho, \ \rho \ \vdash \lnot\lnot\rho$}
        

    \RightLabel{$\lnot_e$}
    \BinaryInfC{$\lnot\lnot\lnot\rho, \ \rho \ \vdash \bot$}

    \RightLabel{$\lnot_i$}
    \UnaryInfC{$\lnot\lnot\lnot\rho \vdash \lnot\rho$}

    \RightLabel{$\Rightarrow_i$}
    \UnaryInfC{$\vdash \lnot\lnot\lnot\rho \Rightarrow \lnot\rho$}
\end{prooftree}

\item {Contraposición: $(\rho \Rightarrow \sigma) \Rightarrow (\lnot\sigma \Rightarrow \lnot\rho)$}

\begin{prooftree}
            \AxiomC{}
            \RightLabel{$_{Ax.}$}
            \UnaryInfC{$\Gamma \vdash \rho \Rightarrow \sigma$}

            
            \AxiomC{}
            \RightLabel{$_{Ax.}$}
            \UnaryInfC{$\Gamma \vdash \rho$}

        \RightLabel{$\Rightarrow_e$}
        \BinaryInfC{$\Gamma = (\rho \Rightarrow \sigma), \ \lnot\sigma, \ \rho \ \vdash \sigma$}
        
        \AxiomC{}
        \RightLabel{$_{Ax.}$}
        \UnaryInfC{$(\rho \Rightarrow \sigma), \ \lnot\sigma, \ \rho \ \vdash \lnot\sigma$}

    \RightLabel{}
    \BinaryInfC{$(\rho \Rightarrow \sigma), \ \lnot\sigma, \ \rho \ \vdash \bot$}

    \RightLabel{$\lnot_i$}
    \UnaryInfC{$(\rho \Rightarrow \sigma), \ \lnot\sigma \ \vdash \lnot\rho$}

    \RightLabel{$\Rightarrow_i$}
    \UnaryInfC{$(\rho \Rightarrow \sigma) \vdash \lnot\sigma \Rightarrow \lnot\rho$}

    \RightLabel{$\Rightarrow_i$}
    \UnaryInfC{$\vdash (\rho \Rightarrow \sigma) \Rightarrow (\lnot\sigma \Rightarrow \lnot\rho)$}
\end{prooftree}

\item {Adjunción: $((\rho \land \sigma) \Rightarrow \tau ) \iff (\rho \Rightarrow \sigma \Rightarrow \tau)$}
\\
\\
Lo separo en dos casos ($\Rightarrow$ y $\Leftarrow$). Primero pruebo uno y después el otro.
\begin{prooftree}
        \AxiomC{}
        \RightLabel{$_{Ax.}$}
        \UnaryInfC{$\Gamma \vdash (\rho \land \sigma) \Rightarrow \tau$}

            \AxiomC{}
            \RightLabel{$_{Ax.}$}
            \UnaryInfC{$\Gamma \vdash \rho$}

            \AxiomC{}
            \RightLabel{$_{Ax.}$}   
            \UnaryInfC{$\Gamma \vdash \sigma$}
            
        \RightLabel{$\land_e$}
        \BinaryInfC{$\Gamma \vdash \rho \land \sigma$}
        
    \RightLabel{$\Rightarrow_e$}
    \BinaryInfC{$\Gamma = ((\rho \land \sigma) \Rightarrow \tau), \ \rho, \ \sigma \ \vdash \tau$}
    
    \RightLabel{$\Rightarrow_i$}
    \UnaryInfC{$((\rho \land \sigma) \Rightarrow \tau), \ \rho \ \vdash \sigma \Rightarrow \tau$}

    \RightLabel{$\Rightarrow_i$}
    \UnaryInfC{$(\rho \land \sigma) \Rightarrow \tau \vdash \rho \Rightarrow \sigma \Rightarrow \tau$}

    \RightLabel{$\Rightarrow_i$}
    \UnaryInfC{$\vdash((\rho \land \sigma) \Rightarrow \tau ) \Rightarrow (\rho \Rightarrow \sigma \Rightarrow \tau)$}
\end{prooftree}
\begin{prooftree}
            \AxiomC{}
            \RightLabel{$_{Ax.}$}
            \UnaryInfC{$\Gamma \vdash \rho \Rightarrow \sigma \Rightarrow \tau$}
            
            \AxiomC{}
            \RightLabel{$_{Ax.}$}
            \UnaryInfC{$\Gamma \vdash \rho \land \sigma$}
            \RightLabel{$\land_{e1}$}
            \UnaryInfC{$\Gamma \vdash \rho$}
            
        \RightLabel{}
        \BinaryInfC{$\Gamma \vdash \sigma \Rightarrow \tau$}

        
        \AxiomC{}
        \RightLabel{$_{Ax.}$}
        \UnaryInfC{$\Gamma \vdash \rho \land \sigma$}
        \RightLabel{$\land_{e2}$}
        \UnaryInfC{$\Gamma \vdash \sigma$}

    \RightLabel{$\Rightarrow_e$}
    \BinaryInfC{$\Gamma = (\rho \Rightarrow \sigma \Rightarrow \tau), \ (\rho \land \sigma) \vdash \tau$}

    \RightLabel{$\Rightarrow_i$}
    \UnaryInfC{$(\rho \Rightarrow \sigma \Rightarrow \tau) \vdash (\rho \land \sigma) \Rightarrow \tau$}

    \RightLabel{$\Rightarrow_i$}
    \UnaryInfC{$\vdash(\rho \Rightarrow \sigma \Rightarrow \tau) \Rightarrow ((\rho \land \sigma) \Rightarrow \tau )$}
\end{prooftree}

\item {de Morgan (I): $\lnot(\rho \lor \sigma) \iff (\lnot\rho \land \lnot\sigma)$}

\item {de Morgan (II): $\lnot(\rho \land \sigma) \iff (\lnot\rho \lor \lnot\sigma)$. Para la
dirección $\Rightarrow$ es necesario usar principios de razo-
namiento clásicos.}

\item {Conmutatividad ($\land$): $(\rho \land \sigma) \Rightarrow (\sigma \land \rho)$}
\begin{prooftree}
        \AxiomC{}
        
        \RightLabel{$_{Ax.}$}
        \UnaryInfC{$\rho \land \sigma \vdash \rho \land \sigma$}
        
        \RightLabel{$\land_{e2}$}
        \UnaryInfC{$\rho \land \sigma \vdash \sigma$}

        \AxiomC{}

        \RightLabel{$_{Ax.}$}
        \UnaryInfC{$\rho \land \sigma \vdash \rho \land \sigma$}
        
        \RightLabel{$\land_{e1}$}
        \UnaryInfC{$\rho \land \sigma \vdash \rho$}

    \RightLabel{$\land_i$}
    \BinaryInfC{$\rho \land \sigma \vdash \sigma \land \rho$}
    
    \RightLabel{$\Rightarrow_i$}
    \UnaryInfC{$\vdash (\rho \land \sigma) \Rightarrow (\sigma \land \rho)$}
\end{prooftree}

\item {Asociatividad ($\land$): $((\rho \land \sigma) \land \tau ) \iff (\rho \land (\sigma \land \tau))$}
\\
\\
Lo separo en dos casos ($\Rightarrow$ y $\Leftarrow$). Primero pruebo uno y después el otro.
\begin{prooftree}
        \AxiomC{}
        
        \RightLabel{$_{Ax.}$}
        \UnaryInfC{$\Gamma \vdash (\rho \land \sigma) \land \tau$}


        \RightLabel{$\land_{e1}$}
        \UnaryInfC{$\Gamma \vdash \rho \land \sigma$}
        
        \RightLabel{$\land_{e1}$}
        \UnaryInfC{$\Gamma \vdash \rho$}
        
            \AxiomC{}
            
            \RightLabel{$_{Ax.}$}
            \UnaryInfC{$\Gamma \vdash (\rho \land \sigma) \land \tau$}

            \RightLabel{$\land_{e1}$}
            \UnaryInfC{$\Gamma \vdash \rho \land \sigma$}
            
            \RightLabel{$\land_{e2}$}
            \UnaryInfC{$\Gamma \vdash \sigma$}

            
            \AxiomC{}
            
            \RightLabel{$_{Ax.}$}
            \UnaryInfC{$\Gamma \vdash (\rho \land \sigma) \land \tau$}
            
            \RightLabel{$\land_{e1}$}
            \UnaryInfC{$\Gamma \vdash \tau$}
            
        \RightLabel{$\land_i$}
        \BinaryInfC{$\Gamma \vdash  \sigma \land \tau$}

    \RightLabel{$\land_i$}
    \BinaryInfC{$\Gamma = (\rho \land \sigma) \land \tau \vdash \rho \land (\sigma \land \tau)$}

    \RightLabel{$\Rightarrow_i$}
    \UnaryInfC{$\vdash ((\rho \land \sigma) \land \tau ) \Rightarrow (\rho \land (\sigma \land \tau))$}
\end{prooftree}

\begin{prooftree}
            \AxiomC{}
            \RightLabel{$_{Ax.}$}
            \UnaryInfC{$\Gamma \vdash \rho \land (\sigma \land \tau)$}
            
            \RightLabel{$\land_{e1}$}
            \UnaryInfC{$\Gamma \vdash \rho$}

            \AxiomC{}
            \RightLabel{$_{Ax.}$}
            \UnaryInfC{{$\Gamma \vdash \rho \land (\sigma \land \tau)$}}

            \RightLabel{$\land_{e2}$}
            \UnaryInfC{$\Gamma \vdash \sigma \land \tau$}

            \RightLabel{$\land_{e1}$}
            \UnaryInfC{$\Gamma \vdash \sigma$}

        \RightLabel{$\land_i$}
        \BinaryInfC{$\Gamma \vdash \rho \land \sigma$}
        
        \AxiomC{}
        \RightLabel{$_{Ax.}$}
        \UnaryInfC{$\Gamma \vdash \rho \land (\sigma \land \tau)$}

        \RightLabel{$\land_{e2}$}
        \UnaryInfC{$\Gamma \vdash \sigma \land \tau$}

        \RightLabel{$\land_{e2}$}
        \UnaryInfC{$\Gamma \vdash \tau$}

    \RightLabel{$\land_i$}
    \BinaryInfC{$\Gamma = \rho \land (\sigma \land \tau) \vdash (\rho \land \sigma) \land \tau$}

    \RightLabel{$\Rightarrow_i$}
    \UnaryInfC{$\vdash (\rho \land (\sigma \land \tau)) \Rightarrow ((\rho \land \sigma) \land \tau )$}
    
\end{prooftree}

\item {Conmutatividad ($\lor$): $(\rho \lor \sigma) \Rightarrow (\sigma \lor \rho)$}

\begin{prooftree}
        \AxiomC{}
        \RightLabel{$_{Ax.}$}
        \UnaryInfC{$(\rho \lor \sigma) \vdash (\rho \lor \sigma)$}

    
        \AxiomC{}
        \RightLabel{$_{Ax.}$}
        \UnaryInfC{$(\rho \lor \sigma), \ \rho \ \vdash \rho$}
        \RightLabel{$\lor_{i2}$}
        \UnaryInfC{$(\rho \lor \sigma), \ \rho \ \vdash (\sigma \lor \rho)$}
        
        \AxiomC{}
        \RightLabel{$_{Ax.}$}
        \UnaryInfC{$(\rho \lor \sigma), \ \sigma \ \vdash \sigma$}
        \RightLabel{$\lor_{i1}$}
        \UnaryInfC{$(\rho \lor \sigma), \ \sigma \ \vdash (\sigma \lor \rho)$}

    \RightLabel{$\lor_e$}
    \TrinaryInfC{$(\rho \lor \sigma) \vdash (\sigma \lor \rho)$}

    \RightLabel{$\Rightarrow_i$}
    \UnaryInfC{$\vdash (\rho \lor \sigma) \Rightarrow (\sigma \lor \rho)$}
\end{prooftree}
\end{enumerate}

\section*{Ejercicio 6}{Demostrar en deducción natural que vale $\vdash \sigma$ para cada una de las siguientes fórmulas. Para estas fórmulas es imprescindible \textbf{usar lógica clásica}:}
\begin{enumerate}[label=\roman*.,font=\itshape]

    \item {Absurdo clásico: $(\lnot\tau \Rightarrow \bot) \Rightarrow \tau$}

    \begin{prooftree}
            \AxiomC{}
            \RightLabel{$_{Ax.}$}
            \UnaryInfC{$(\lnot\tau \Rightarrow \bot), \lnot\tau \vdash \lnot\tau \Rightarrow \bot$}
            
            \AxiomC{}
            \RightLabel{$_{Ax.}$}
            \UnaryInfC{$(\lnot\tau \Rightarrow \bot), \lnot\tau \vdash \lnot\tau$}

        \RightLabel{$\Rightarrow_e$}
        \BinaryInfC{$(\lnot\tau \Rightarrow \bot), \lnot\tau \vdash \bot$}

        \RightLabel{$_{PBC}$}
        \UnaryInfC{$(\lnot\tau \Rightarrow \bot) \vdash \tau$}        

        \RightLabel{$\Rightarrow_i$}
        \UnaryInfC{$\vdash (\lnot\tau \Rightarrow \bot) \Rightarrow \tau$}
    \end{prooftree}

\end{enumerate}

\end{document}
